\section{Motivation}

Sokoban occupies a distinctive position in artificial intelligence research. Unlike games such as chess or Go, where positions can be evaluated and mistakes can often be recovered from, Sokoban is a puzzle game where a wrong move can prevent the player from proceeding, often forcing a full reset. For example, a box can be permanently trapped in a corner or the player can be trapped by boxes. The irreversibility and the game's formal properties makes this game both NP-hard \cite{fryers1995sokoban} and PSPACE-complete \cite{culberson1997sokoban} placing it at the upper limits of complexity for discrete planning.

A* search guided by a learned cost-to-go heuristic is one of the most effective general approaches to Sokoban. The heuristic estimates the number of steps remaining to the solution from any board state, and its accuracy directly determines how efficiently the search tree is explored. The dominant paradigm for learning such heuristics is to fix a neural network architecture --- typically a multi-layer convolutional or fully connected network --- and train its weights by gradient descent on supervised labels derived from solution trajectories. This approach is effective but has a significant drawback: gradient descent can only optimize within the fixed topology chosen before training. If that topology is overparameterized for the problem, the resulting heuristic wastes memory and computation; if it lacks the right inductive structure, no amount of training can compensate.

This thesis investigates an alternative: using neuroevolution to search over both the topology and the weights of the heuristic simultaneously. Rather than committing to an architecture in advance, evolutionary methods can discover the network structure that best fits the heuristic learning task, potentially finding solutions that are far more compact than those produced by gradient descent over a fixed architecture. The question motivating this work is whether a neuroevolved heuristic can match the accuracy of a hand-designed convolutional baseline while using orders of magnitude fewer parameters.

\section{Problem Statement}

Sokoban presents two fundamental challenges that set it apart from most reinforcement learning benchmarks.

The first is irreversibility. Every action in Sokoban is potentially permanent: once a box is pushed against a wall or into a corner that is not a goal position, it cannot be retrieved. This creates deadlock states from which the only escape is a full puzzle reset. A successful agent must therefore evaluate not only the immediate effect of each action but its consequences many steps into the future, making shallow strategies inadequate.

The second challenge is sparse rewards. The agent receives no feedback signal until all boxes have been placed on their goal positions. In a $10 \times 10$ grid with multiple boxes and goals, reaching the solution purely by exploration is statistically unlikely, leaving standard reinforcement learning methods with almost no gradient signal to learn from in early training.

Beyond these, the design of the heuristic function itself presents a practical challenge. A fixed convolutional architecture trained by gradient descent carries a parameter count determined by human design choices, not by the complexity of the task. The CNN baseline used in this work has 2,031,809 parameters, the majority of which reside in a single dense layer connecting the convolutional feature map to the hidden state. Whether all of these parameters contribute meaningfully to heuristic quality, or whether a much smaller network discovered by evolutionary search can achieve comparable accuracy, is the central empirical question of this thesis.

\section{Contributions}

This thesis makes three primary contributions.

First, we implement a CNN baseline heuristic --- a fixed four-layer convolutional network with 2,031,809 parameters trained by gradient descent with an asymmetric Huber loss --- and establish its cost-to-go prediction accuracy and A* solve rate on the Boxoban benchmark as the reference point for comparison.

Second, we apply NEAT to evolve a feed-forward heuristic network from a minimal starting topology of approximately 26 parameters. NEAT searches simultaneously over network topology and weights, adding nodes and connections only where doing so improves the negative-MAE fitness signal, and uses speciation to protect structural innovations from premature elimination. The evolved heuristic is expected to be several orders of magnitude smaller than the CNN baseline.

Third, we apply CoDeepNEAT to coevolve populations of dense-dropout modules and blueprint topologies, with each assembled genome trained by gradient descent for a small number of epochs per generation. This hybrid approach combines the architectural flexibility of evolutionary search with the weight-optimization efficiency of gradient descent, producing heuristics in the range of 32,000 to 170,000 parameters --- between 12$\times$ and 63$\times$ smaller than the CNN baseline.

All three heuristics are trained on cost-to-go labels derived from Breadth-First Search solution trajectories over the Boxoban dataset and evaluated on held-out levels by mean absolute error and Weighted A* solve rate.

\begin{figure}[h]
    \centering
    \includegraphics[width=0.8\textwidth]{sokoban}
    \caption{Screenshot of Sokoban as seen on https://www.sokoban.jp/. The player character is wearing a yellow hardhat and pushing a box into a wall while surrounded by brick walls. The box needs to be placed on the red dot for solution.}
    \label{fig:sokoban}
\end{figure}
